\documentclass[12pt]{article}

% Idioma y codificación
\usepackage[spanish, es-tabla, es-notilde]{babel}       %es-tabla para que se titule "Tabla"
\usepackage[utf8]{inputenc}

% Márgenes
\usepackage[a4paper,top=3cm,bottom=2.5cm,left=3cm,right=3cm]{geometry}

% Comentarios de bloque
\usepackage{verbatim}

% Paquetes de links
\usepackage[hidelinks]{hyperref}    % Permite enlaces
\usepackage{url}                    % redirecciona a la web

% Más opciones para enumeraciones
\usepackage{enumitem}

% Personalizar la portada
\usepackage{titling}

% Paquetes de tablas
\usepackage{multirow}

% Para añadir el símbolo de euro
\usepackage{eurosym}

% Paquete para subrayar con diferentes estilos
\usepackage[normalem]{ulem}


%------------------------------------------------------------------------

%Paquetes de figuras
\usepackage{caption}
\usepackage{subcaption} % Figuras al lado de otras
\usepackage{float}      % Poner figuras en el sitio indicado H.


% Paquetes de imágenes
\usepackage{graphicx}       % Paquete para añadir imágenes
\usepackage{transparent}    % Para manejar la opacidad de las figuras

% Paquete para usar colores
\usepackage[dvipsnames, table, xcdraw]{xcolor}
\usepackage{pagecolor}      % Para cambiar el color de la página

% Habilita tamaños de fuente mayores
\usepackage{fix-cm}

% Para los gráficos
\usepackage{tikz}
\usepackage{forest}

% Para poder situar los nodos en los grafos
\usetikzlibrary{positioning}

% Para evitar problemas entre babel y tikz
\usetikzlibrary{babel}


%------------------------------------------------------------------------

% Paquetes de matemáticas
\usepackage{mathtools, amsfonts, amssymb, mathrsfs}
\usepackage[makeroom]{cancel}     % Simplificar tachando
\usepackage{polynom}    % Divisiones y Ruffini
\usepackage{units} % Para poner fracciones diagonales con \nicefrac

\usepackage{pgfplots}   %Representar funciones
\pgfplotsset{compat=1.18}  % Versión 1.18

\usepackage{tikz-cd}    % Para usar diagramas de composiciones
\usetikzlibrary{calc}   % Para usar cálculo de coordenadas en tikz

%Definición de teoremas, etc.
\usepackage{amsthm}
%\swapnumbers   % Intercambia la posición del texto y de la numeración

\theoremstyle{plain}

\makeatletter
\@ifclassloaded{article}{
  % --- Configuración para ARTICLE ---
  \newtheorem{teo}{Teorema}[section]
  \theoremstyle{definition}
  \newtheorem{definicion}{Definición}[section]
  \newtheorem{ejercicio}{Ejercicio}[section]

  \renewcommand{\theteo}{%
    \ifnum\value{section}=0 \arabic{teo}\else \thesection.\arabic{teo}\fi
  }
  \renewcommand{\thedefinicion}{%
    \ifnum\value{section}=0 \arabic{definicion}\else \thesection.\arabic{definicion}\fi
  }
  \renewcommand{\theejercicio}{%
    \ifnum\value{section}=0 \arabic{ejercicio}\else \thesection.\arabic{ejercicio}\fi
  }
}{
  % --- Configuración para BOOK / REPORT ---
  \newtheorem{teo}{Teorema}[chapter]
  \theoremstyle{definition}
  \newtheorem{definicion}{Definición}[chapter]
  \newtheorem{ejercicio}{Ejercicio}[section]

  % Si no hay sección: Capítulo.Teorema (ej. 2.1)
  % Si hay sección:    Capítulo.Sección.Teorema (ej. 2.3.1)
  \renewcommand{\theteo}{%
    \ifnum\value{section}=0
      \thechapter.\arabic{teo}%
    \else
      \thesection.\arabic{teo}% (\thesection ya contiene el número de capítulo)
    \fi
  }
  \renewcommand{\thedefinicion}{%
    \ifnum\value{section}=0
      \thechapter.\arabic{definicion}%
    \else
      \thesection.\arabic{definicion}%
    \fi
  }
  \renewcommand{\theejercicio}{%
    \ifnum\value{section}=0
      \arabic{ejercicio}%
    \else
      \thesection.\arabic{ejercicio}% Formato por defecto: ej. 2.2.1
    \fi
  }
}
\makeatother

% Teoremas dependientes o con contador compartido
\newtheorem{coro}{Corolario}[teo]
\newtheorem{prop}[teo]{Proposición}
\newtheorem{lema}[teo]{Lema}
\newtheorem*{lema*}{Lema}

\theoremstyle{remark}
\newtheorem*{observacion}{Observación}

\theoremstyle{definition}
\newtheorem*{notacion}{Notación}
\newtheorem*{ejemplo}{Ejemplo}
\newtheorem*{ejercicio*}{Ejercicio}


% \renewcommand\qedsymbol{$\blacksquare$}         % Cambiar símbolo QED
%------------------------------------------------------------------------

% Paquetes para encabezados
\usepackage{fancyhdr}
\pagestyle{fancy}
\fancyhf{}

\newcommand{\helv}{ % Modificación tamaño de letra
\fontfamily{}\fontsize{12}{12}\selectfont}
\setlength{\headheight}{15pt} % Amplía el tamaño del índice


%\usepackage{lastpage}   % Referenciar última pag   \pageref{LastPage}
%\fancyfoot[C]{%
%  \begin{minipage}{\textwidth}
%    \centering
%    ~\\
%    \thepage\\
%    \href{https://losdeldgiim.github.io/}{\texttt{\footnotesize losdeldgiim.github.io}}
%  \end{minipage}
%}
\fancyfoot[C]{\thepage}
\fancyfoot[R]{\href{https://losdeldgiim.github.io/}{\texttt{\footnotesize losdeldgiim.github.io}}}

%------------------------------------------------------------------------

% Conseguir que no ponga "Capítulo 1". Sino solo "1."
\makeatletter
\@ifclassloaded{book}{
  \renewcommand{\chaptermark}[1]{\markboth{\thechapter.\ #1}{}} % En el encabezado
    
  \renewcommand{\@makechapterhead}[1]{%
  \vspace*{50\p@}%
  {\parindent \z@ \raggedright \normalfont
    \ifnum \c@secnumdepth >\m@ne
      \huge\bfseries \thechapter.\hspace{1em}\ignorespaces
    \fi
    \interlinepenalty\@M
    \Huge \bfseries #1\par\nobreak
    \vskip 40\p@
  }}
}
\makeatother

%------------------------------------------------------------------------
% Paquetes de cógido
\usepackage{minted}
\renewcommand\listingscaption{Código fuente}

\usepackage{fancyvrb}
% Personaliza el tamaño de los números de línea
\renewcommand{\theFancyVerbLine}{\small\arabic{FancyVerbLine}}

% Estilo para C++
\newminted{cpp}{
    frame=lines,
    framesep=2mm,
    baselinestretch=1.2,
    linenos,
    escapeinside=||
}

% para minted
\definecolor{LightGray}{rgb}{0.95,0.95,0.92}
\setminted{
    linenos=true,
    stepnumber=5,
    numberfirstline=true,
    autogobble,
    breaklines=true,
    breakautoindent=true,
    breaksymbolleft=,
    breaksymbolright=,
    breaksymbolindentleft=0pt,
    breaksymbolindentright=0pt,
    breaksymbolsepleft=0pt,
    breaksymbolsepright=0pt,
    fontsize=\footnotesize,
    bgcolor=LightGray,
    numbersep=10pt
}


\usepackage{listings} % Para incluir código desde un archivo

\renewcommand\lstlistingname{Código Fuente}
\renewcommand\lstlistlistingname{Índice de Códigos Fuente}

% Definir colores
\definecolor{vscodepurple}{rgb}{0.5,0,0.5}
\definecolor{vscodeblue}{rgb}{0,0,0.8}
\definecolor{vscodegreen}{rgb}{0,0.5,0}
\definecolor{vscodegray}{rgb}{0.5,0.5,0.5}
\definecolor{vscodebackground}{rgb}{0.97,0.97,0.97}
\definecolor{vscodelightgray}{rgb}{0.9,0.9,0.9}

% Configuración para el estilo de C similar a VSCode
\lstdefinestyle{vscode_C}{
  backgroundcolor=\color{vscodebackground},
  commentstyle=\color{vscodegreen},
  keywordstyle=\color{vscodeblue},
  numberstyle=\tiny\color{vscodegray},
  stringstyle=\color{vscodepurple},
  basicstyle=\scriptsize\ttfamily,
  breakatwhitespace=false,
  breaklines=true,
  captionpos=b,
  keepspaces=true,
  numbers=left,
  numbersep=5pt,
  showspaces=false,
  showstringspaces=false,
  showtabs=false,
  tabsize=2,
  frame=tb,
  framerule=0pt,
  aboveskip=10pt,
  belowskip=10pt,
  xleftmargin=10pt,
  xrightmargin=10pt,
  framexleftmargin=10pt,
  framexrightmargin=10pt,
  framesep=0pt,
  rulecolor=\color{vscodelightgray},
  backgroundcolor=\color{vscodebackground},
}

%------------------------------------------------------------------------
% La etiqueda IA sirve para indicar a partir de donde le pasamos el preambulo a la IA para el escaneo
% IA

% Comandos definidos
\newcommand{\bb}[1]{\mathbb{#1}}
\newcommand{\cc}[1]{\mathcal{#1}}

% I prefer the slanted \leq
\let\oldleq\leq % save them in case they're every wanted
\let\oldgeq\geq
\renewcommand{\leq}{\leqslant}
\renewcommand{\geq}{\geqslant}

% Si y solo si
\newcommand{\sii}{\iff}

% MCD y MCM
\DeclareMathOperator{\mcd}{mcd}
\DeclareMathOperator{\mcm}{mcm}

% Signo
\DeclareMathOperator{\sgn}{sgn}

% Imagen
\DeclareMathOperator{\Ima}{Im}

% Aplicación identidad
\newcommand{\id}{\mathrm{Id}}

% Letras griegas
\newcommand{\eps}{\epsilon}
\newcommand{\veps}{\varepsilon}
\newcommand{\lm}{\lambda}

\newcommand{\ol}{\overline}
\newcommand{\ul}{\underline}
\newcommand{\wt}{\widetilde}
\newcommand{\wh}{\widehat}

\let\oldvec\vec
\renewcommand{\vec}{\overrightarrow}

% Derivadas parciales
\newcommand{\del}[2]{\frac{\partial #1}{\partial #2}}
\newcommand{\Del}[3]{\frac{\partial^{#1} #2}{\partial #3^{#1}}}
\newcommand{\deld}[2]{\dfrac{\partial #1}{\partial #2}}
\newcommand{\Deld}[3]{\dfrac{\partial^{#1} #2}{\partial #3^{#1}}}


\newcommand{\AstIg}{\stackrel{(\ast)}{=}}
\newcommand{\Hop}{\stackrel{L'H\hat{o}pital}{=}}

\newcommand{\red}[1]{{\color{red}#1}} % Para integrales, destacar los cambios.

% Método de integración
\newcommand{\MetInt}[2]{
    \left[\begin{array}{c}
        #1 \\ #2
    \end{array}\right]
}

% Declarar aplicaciones
% 1. Nombre aplicación
% 2. Dominio
% 3. Codominio
% 4. Variable
% 5. Imagen de la variable
\newcommand{\Func}[5]{
    \begin{equation*}
        \begin{array}{rrll}
            \displaystyle #1:& \displaystyle  #2 & \longrightarrow & \displaystyle  #3\\
               & \displaystyle  #4 & \longmapsto & \displaystyle  #5
        \end{array}
    \end{equation*}
}

%------------------------------------------------------------------------


\newcommand{\myhref}[2]{\href{#1}{\color{blue}\uline{#2}}}

% // TODO: Publicar
\begin{document}

    % 1. Foto de fondo
    % 2. Título
    % 3. Encabezado Izquierdo
    % 4. Color de fondo
    % 5. Coord x del titulo
    % 6. Coord y del titulo
    % 7. Fecha

    
    % 1. Foto de fondo
% 2. Título
% 3. Encabezado Izquierdo
% 4. Color de fondo
% 5. Coord x del titulo
% 6. Coord y del titulo
% 7. Fecha
% 8. Autor

\newcommand{\portada}[8]{

    \portadaBase{#1}{#2}{#3}{#4}{#5}{#6}{#7}{#8}
    \portadaBook{#1}{#2}{#3}{#4}{#5}{#6}{#7}{#8}
}

\newcommand{\portadaExamen}[8]{

    \portadaBase{#1}{#2}{#3}{#4}{#5}{#6}{#7}{#8}
    \portadaArticle{#1}{#2}{#3}{#4}{#5}{#6}{#7}{#8}
}




\newcommand{\portadaBase}[8]{

    % Tiene la portada principal y la licencia Creative Commons
    
    % 1. Foto de fondo
    % 2. Título
    % 3. Encabezado Izquierdo
    % 4. Color de fondo
    % 5. Coord x del titulo
    % 6. Coord y del titulo
    % 7. Fecha
    % 8. Autor
    
    
    \thispagestyle{empty}               % Sin encabezado ni pie de página
    \newgeometry{margin=0cm}        % Márgenes nulos para la primera página
    
    
    % Encabezado
    \fancyhead[L]{\helv #3}
    \fancyhead[R]{\helv \nouppercase{\leftmark}}
    
    
    \pagecolor{#4}        % Color de fondo para la portada
    
    \begin{figure}[p]
        \centering
        \transparent{0.3}           % Opacidad del 30% para la imagen
        
        \includegraphics[width=\paperwidth, keepaspectratio]{../../_assets/#1}
    
        \begin{tikzpicture}[remember picture, overlay]
            \node[anchor=north west, text=white, opacity=1, font=\fontsize{60}{90}\selectfont\bfseries\sffamily, align=left] at (#5, #6) {#2};
            
            \node[anchor=south east, text=white, opacity=1, font=\fontsize{12}{18}\selectfont\sffamily, align=right] at (9.7, 3) {\href{https://losdeldgiim.github.io/}{\textbf{Los Del DGIIM}, \texttt{\footnotesize losdeldgiim.github.io}}};
            
            \node[anchor=south east, text=white, opacity=1, font=\fontsize{12}{15}\selectfont\sffamily, align=right] at (9.7, 1.8) {Doble Grado en Ingeniería Informática y Matemáticas\\Universidad de Granada};
        \end{tikzpicture}
    \end{figure}
    
    
    \restoregeometry        % Restaurar márgenes normales para las páginas subsiguientes
    \pagecolor{white}       % Restaurar el color de página
    
    
    \newpage
    \thispagestyle{empty}               % Sin encabezado ni pie de página
    \begin{tikzpicture}[remember picture, overlay]
        \node[anchor=south west, inner sep=3cm] at (current page.south west) {
            \begin{minipage}{0.5\paperwidth}
                \href{https://creativecommons.org/licenses/by-nc-nd/4.0/}{
                    \includegraphics[height=2cm]{../../_assets/Licencia.png}
                }\vspace{1cm}\\
                Esta obra está bajo una
                \href{https://creativecommons.org/licenses/by-nc-nd/4.0/}{
                    Licencia Creative Commons Atribución-NoComercial-SinDerivadas 4.0 Internacional (CC BY-NC-ND 4.0).
                }\\
    
                Eres libre de compartir y redistribuir el contenido de esta obra en cualquier medio o formato, siempre y cuando des el crédito adecuado a los autores originales y no persigas fines comerciales. 
            \end{minipage}
        };
    \end{tikzpicture}
    
    
    
    % 1. Foto de fondo
    % 2. Título
    % 3. Encabezado Izquierdo
    % 4. Color de fondo
    % 5. Coord x del titulo
    % 6. Coord y del titulo
    % 7. Fecha
    % 8. Autor


}


\newcommand{\portadaBook}[8]{

    % 1. Foto de fondo
    % 2. Título
    % 3. Encabezado Izquierdo
    % 4. Color de fondo
    % 5. Coord x del titulo
    % 6. Coord y del titulo
    % 7. Fecha
    % 8. Autor

    % Personaliza el formato del título
    \pretitle{\begin{center}\bfseries\fontsize{42}{56}\selectfont}
    \posttitle{\par\end{center}\vspace{2em}}
    
    % Personaliza el formato del autor
    \preauthor{\begin{center}\Large}
    \postauthor{\par\end{center}\vfill}
    
    % Personaliza el formato de la fecha
    \predate{\begin{center}\huge}
    \postdate{\par\end{center}\vspace{2em}}
    
    \title{#2}
    \author{\href{https://losdeldgiim.github.io/}{Los Del DGIIM, \texttt{\large losdeldgiim.github.io}}
    \\ \vspace{0.5cm}#8}
    \date{Granada, #7}
    \maketitle
    
    \tableofcontents
}




\newcommand{\portadaArticle}[8]{

    % 1. Foto de fondo
    % 2. Título
    % 3. Encabezado Izquierdo
    % 4. Color de fondo
    % 5. Coord x del titulo
    % 6. Coord y del titulo
    % 7. Fecha
    % 8. Autor

    % Personaliza el formato del título
    \pretitle{\begin{center}\bfseries\fontsize{42}{56}\selectfont}
    \posttitle{\par\end{center}\vspace{2em}}
    
    % Personaliza el formato del autor
    \preauthor{\begin{center}\Large}
    \postauthor{\par\end{center}\vspace{3em}}
    
    % Personaliza el formato de la fecha
    \predate{\begin{center}\huge}
    \postdate{\par\end{center}\vspace{5em}}
    
    \title{#2}
    \author{\href{https://losdeldgiim.github.io/}{Los Del DGIIM, \texttt{\large losdeldgiim.github.io}}
    \\ \vspace{0.5cm}#8}
    \date{Granada, #7}
    \thispagestyle{empty}               % Sin encabezado ni pie de página
    \maketitle
    \vfill
}
    \portadaExamen{ffccA4.jpg}{Análisis Funcional\\Examen XXI}{Análisis Funcional. Examen XXI}{MidnightBlue}{-8}{28}{2026}{}

    \begin{description}
        \item[Asignatura] Análisis Funcional.
        \item[Curso Académico] 2026/2027.
        \item[Grado] Doble Grado en Ingeniería Informática y Matemáticas.
        \item[Grupo] Único.
        \item[Profesor] Miguel Martín Suárez.
        \item[Descripción] Parcial 2.
        \item[Fecha] 25 de noviembre de 2026.
        \item[Duración] 50 minutos.
    
    \end{description}
    \newpage


    % ------------------------------------

En el examen había 7 preguntas tipo test, cada una puntuando $1.5$ puntos (por lo que elexamen estaba puntuado sobre $10.5$ puntos). Cada pregunta tenía 4 respuestas posibles, de las cuales una, dos, tres o cuatro podían ser correctas. La puntuación de cada pregunta se calculaba de la siguiente manera:
\begin{itemize}
    \item Si se selecciona alguna respuesta incorrecta, la puntuación de la pregunta es 0.
    \item Si se seleccionan $k$ respuestas correctas y no se selecciona ninguna incorrecta, la puntuación de la pregunta es:
    \begin{equation*}
        \frac{k}{n} \cdot 1.5
    \end{equation*}
    donde $n$ es el número total de respuestas correctas de la pregunta.
\end{itemize}


Aunque cada alumno tenía un examen distinto que incluía 7 preguntas tipo test, se incluyen todas las preguntas para que el estudiantado tenga una mayor variedad de ejercicios para practicar.
\begin{ejercicio}
    Sea $X$ un subespacio cerrado de $\ell_{2}$ de dimensión infinita ¿Puede ser $X$ isomorfo a $\ell_1$?
    \begin{enumerate}[label={\arabic*)}]
        \item Sí, porque $\ell_1 \subseteq \ell_{2}$
        \item No, porque $\ell_1 \nsubseteq \ell_{2}$
        \item No, porque $\ell_1$ es separable y $X$ no lo es
        \item No, porque $X^*$ es separable y $\ell_1^*$ no lo es
    \end{enumerate}
\end{ejercicio}

\begin{ejercicio}
    Sea $\psi \colon \ell_2 \to \bb{R}$ el funcional lineal dado por $\psi(x) = \sum\limits_{n=1}^\infty \frac{x(n)}{n}$ para $x \in \ell_2$:
    \begin{enumerate}[label={\arabic*)}]
        \item $\psi$ alcanza su norma porque $\ell_2$ es reflexivo
        \item $\psi$ alcanza su norma porque $\ell_2$ tiene dimensión finita
        \item $\psi$ no alcanza su norma
        \item $\psi$ no alcanza su norma porque $\ell_2$ no es completo
    \end{enumerate}
\end{ejercicio}

\begin{ejercicio}
    En un espacio normado $X$ tomamos un conjunto convexo y cerrado $C$ y un punto $x_0 \notin C$:
    \begin{enumerate}[label={\arabic*)}]
        \item existe $f \in X^*$ tal que $\sup \Re f(C) \leq \Re f(x_0)$
        \item existe $f \in X^*$ tal que $\sup \Re f(C) < \Re f(x_0)$
        \item no se puede asegurar que $C$ se pueda separar de $x_0$ si el interior de $C$ es vacío
        \item existe un único $f \in X^*$ tal que $\sup \Re f(C) \leq \Re f(x_0)$
    \end{enumerate}
\end{ejercicio}

\begin{ejercicio}
    Sea $X=c_0$. Entonces:
    \begin{enumerate}[label={\arabic*)}]
        \item $X^{***}$ contiene un subespacio isométricamente isomorfo a $\ell_1$
        \item $X^{**}$ contiene un subespacio isomorfo a $\ell_1$
        \item Contiene un cociente isomorfo a $\ell_1$
        \item Contiene un subespacio isomorfo a $\ell_1$ ???
    \end{enumerate}
\end{ejercicio}

\begin{ejercicio}
    Sea $X$ un espacio normado. Entonces:
    \begin{enumerate}[label={\arabic*)}]
        \item $X$ es isométricamente isomorfo a un subespacio cerrado de $X^{**}$.
        \item $X$ es isométricamente isomorfo a un subespacio de $X^{**}$.
        \item $X^*$ es completo.
        \item Si $X^*$ es separable, entonces todo subespacio de $X$ es separable.
    \end{enumerate}
\end{ejercicio}

\begin{ejercicio}
    Sea $X$ un espacio de Banach reflexivo y separable. Entonces:
    \begin{enumerate}[label={\arabic*)}]
        \item Si $M$ es un subespacio cerrado de $X$, entonces $M$ es proximinal.
        \item $X^*$ es reflexivo.
        \item $X^*$ es separable.
        \item Tiene un subespacio isométricamente isomorfo a $\ell_1$.
    \end{enumerate}
\end{ejercicio}

\begin{ejercicio}
    Dados $X,Y,Z$ espacios normados y $T\in \cc{L}(X,Y)$ entonces:
    \begin{enumerate}[label={\arabic*)}]
        \item El operador $T^*$ cumple $T^* \in \cc{L}(Y^*,X^*)$ y $\|T^*\| = \|T\|$.
        \item Dados $S\in\in \cc{L}(X,Y)$ y $T\in \cc{L}(Y,Z)$, entonces $(S\circ T)^* = S^*\circ T^*$.
        \item Si $T$ es un isomorfismo, entonces $T^*$ es un isomorfismo y $(T^{-1})^* = (T^*)^{-1}$.
        \item Si $J_X\colon X\to X^{**}$ y $J_Y\colon Y\to Y^{**}$ son las inyecciones canónicas, entonces $T^{**}\circ J_X = J_Y\circ T$.
    \end{enumerate}
\end{ejercicio}

\begin{ejercicio}
    Sea $X$ un espacio de Banach reflexivo.
    \begin{enumerate}[label={\arabic*)}]
        \item Si $T*\colon X^*\to X^*$ es un isomorfismo, entonces $T\colon X\to X$ es un isomorfismo.
    \end{enumerate}
\end{ejercicio}


    \newpage
    \setcounter{ejercicio}{0} % Reiniciar contador de ejercicios
    \noindent
    \textbf{Solución.}

\begin{ejercicio}
    Sea $X$ un subespacio cerrado de $\ell_{2}$ de dimensión infinita ¿Puede ser $X$ isomorfo a $\ell_1$?
    \begin{enumerate}[label={\arabic*)}]
        \item Sí, porque $\ell_1 \subseteq \ell_{2}$
        \item No, porque $\ell_1 \nsubseteq \ell_{2}$
        \item No, porque $\ell_1$ es separable y $X$ no lo es
        \item No, porque $X^*$ es separable y $\ell_1^*$ no lo es
    \end{enumerate}

    Las dos primeras opciones no están relacionadas, por lo que no son ciertas. Veamos en primer lugar una forma de ver que no pueden ser isomorfos.
    \begin{itemize}
        \item Como $X$ es un subespacio cerrado de $\ell_{2}$, que es reflexivo, se deduce que $X$ es reflexivo. Como $\ell_1$ no es reflexivo, se deduce que $X$ y $\ell_1$ no pueden ser isomorfos.
    \end{itemize}

    No obstante, esto no ayuda a resolver este apartado. Tenemos que:
    \begin{itemize}
        \item $\ell_1$ es separable.
        \item $\ell_{1}^*$ es isométricamente isomorfo a $\ell_\infty$, que no es separable.
        \item $X^*$ es isométricamente isomorfo a $\ell_{2}^*/X^\perp$. Como $\ell_{2}^*$ es separable y la aplicación proyección al cociente es continua y sobreyectiva, se deduce que $\ell_{2}^*/X^\perp$ es separable. Por lo tanto, $X^*$ es separable.
        \item Como $\ell_{2}$ es separable, entonces $X$ es separable.
    \end{itemize}

    Por tanto, por descarte, la única afirmación verdadera es la \textbf{4)}. Esta además es verdadera pues, si $X$ y $\ell_1$ fuesen isomorfos, entonces $X^*$ y $\ell_1^*$ serían isomorfos, lo que es imposible pues $X^*$ es separable y $\ell_1^*$ no lo es.
\end{ejercicio}

\begin{ejercicio}
    Sea $\psi \colon \ell_2 \to \bb{R}$ el funcional lineal dado por $\psi(x) = \sum\limits_{n=1}^\infty \frac{x(n)}{n}$ para $x \in \ell_2$:
    \begin{enumerate}[label={\arabic*)}]
        \item $\psi$ alcanza su norma porque $\ell_2$ es reflexivo
        \item $\psi$ alcanza su norma porque $\ell_2$ tiene dimensión finita
        \item $\psi$ no alcanza su norma
        \item $\psi$ no alcanza su norma porque $\ell_2$ no es completo
    \end{enumerate}

    Como $\psi\in \ell_2^*$ y $\ell_2$ es reflexivo, se deduce que $\psi$ alcanza su norma (equivalentemente, $|\psi|$ alcanza su supremo en la bola unidad de $\ell_2$). Por tanto, opción \textbf{1)} es verdadera.
\end{ejercicio}

\begin{ejercicio}
    En un espacio normado $X$ tomamos un conjunto convexo y cerrado $C$ y un punto $x_0 \notin C$:
    \begin{enumerate}[label={\arabic*)}]
        \item existe $f \in X^*$ tal que $\sup \Re f(C) \leq \Re f(x_0)$
        \item existe $f \in X^*$ tal que $\sup \Re f(C) < \Re f(x_0)$
        \item no se puede asegurar que $C$ se pueda separar de $x_0$ si el interior de $C$ es vacío
        \item existe un único $f \in X^*$ tal que $\sup \Re f(C) \leq \Re f(x_0)$
    \end{enumerate}

    Suponemos que $C\neq \emptyset$, pues si $C=\emptyset$ no se puede considerar su separación de cualquier otro conjunto. En este caso, puesto que ambos conjuntos podrían tener interior vacío, no se puede aplicar de manera directa la separación en espacios normados. No obstante, puesto que $d(x_0, C) > 0$, se puede aplicar la separación fuerte en espacios normados. Como $\{x_0\},C$ son conjuntos no vacíos y convexos tal que $d(\{x_0\}, C) = d(x_0, C) > 0$, se deduce que existen $f \in X^*$ y $\alpha \in \bb{R}$ tales que:
    \begin{equation*}
        \sup \Re f(C) \leq  \alpha<\alpha + d(x_0, C) \leq \Re f(x_0).
    \end{equation*}

    Por lo tanto, las opciones \textbf{1)} y \textbf{2)} son verdaderas, y la opción \textbf{3)} es falsa. La opción \textbf{4)} es falsa porque no hay unicidad del funcional separador (ni siquiera exigiendo que tenga norma $1$). 

    Por ejemplo, en $X = \mathbb{R}^2$ con la norma euclídea, consideremos el conjunto convexo y cerrado $C = \{(0,0)\}$ y el punto $x_0 = (1,0) \notin C$. Los siguientes dos funcionales distintos $f_1, f_2 \in (\mathbb{R}^2)^*$, ambos de norma $1$, separan $C$ de $x_0$:
    \begin{align*}
        f_1(x,y) &= x \quad \Longrightarrow \quad \sup f_1(C) = 0 < 1 = f_1(x_0), \\
        f_2(x,y) &= \frac{1}{\sqrt{2}}(x+y) \quad \Longrightarrow \quad \sup f_2(C) = 0 < \frac{1}{\sqrt{2}} = f_2(x_0).
    \end{align*}
\end{ejercicio}


\end{document}